\chapter{Inverse Trigonometric Functions}
\label{ch:ch07-inverse-functions}

\section{Why Ordinary Inverses Fail}

The previous chapter closed on a difficulty worth stating plainly: the
equation $\sin\theta = \tfrac12$ does not have a single solution, since
$\theta = \pi/6, 5\pi/6, 13\pi/6, 17\pi/6, \dots$ all satisfy it. If several
different inputs produce the same output, no inverse function can determine
which input was originally intended, and sine, as a result, cannot be
reversed uniquely across its full domain. The remedy is simple in
principle: rather than working with the entire graph of sine, a portion of
it is chosen on which every output occurs exactly once, and the inverse is
defined only on that restricted portion.

\section{One-to-One Functions and Invertibility}

\begin{definition}
A function $f$ is \emph{one-to-one} if different inputs always produce
different outputs; equivalently, $f(a) = f(b)$ implies $a = b$. A function
has an inverse if and only if it is one-to-one on its domain.
\end{definition}

Sine fails to be one-to-one on all of $\mathbb{R}$, since for instance
$\sin(\pi/6) = \sin(5\pi/6) = \tfrac12$, and consequently sine has no
inverse on its full domain. If the domain is restricted to
$\left[-\pi/2, \pi/2\right]$, however, every output value between $-1$ and
$1$ occurs exactly once, and on that interval sine becomes invertible.

\section{Restricting the Domain of Sine}

\begin{theorem}
The sine function is one-to-one on $\left[-\pi/2, \pi/2\right]$.
\end{theorem}

\begin{proof}
As $\theta$ moves from $-\pi/2$ to $\pi/2$, the corresponding point on the
unit circle moves continuously from $(0,-1)$ to $(0,1)$, and its vertical
coordinate increases strictly throughout this motion. Since sine measures
the vertical coordinate, distinct angles in this interval produce distinct
sine values, so sine is one-to-one there.
\end{proof}

\section{The Arcsine Function}

Because sine is one-to-one on this restricted domain, it has an inverse
defined there.

\begin{definition}
The \emph{arcsine} function is defined by $\arcsin x = \theta$ if and only
if $\sin\theta = x$ and $-\pi/2 \le \theta \le \pi/2$. Its domain is
$[-1,1]$ and its range is $\left[-\pi/2, \pi/2\right]$.
\end{definition}

Thus $\arcsin(0) = 0$, $\arcsin(1) = \pi/2$, and $\arcsin(1/2) = \pi/6$.
The last of these deserves emphasis: although many angles have sine equal
to $1/2$, arcsine always returns the unique one of them lying in its
restricted range.

\section{Arccosine and Arctangent}

The same idea extends to cosine and tangent, each restricted to whichever
interval makes it one-to-one.

\begin{definition}
Cosine is one-to-one on $[0,\pi]$, and $\arccos x = \theta$ means
$\cos\theta = x$ with $0 \le \theta \le \pi$.
\end{definition}

Hence $\arccos(1) = 0$, $\arccos(0) = \pi/2$, and $\arccos(-1) = \pi$.

\begin{definition}
Tangent is one-to-one on $\left(-\pi/2, \pi/2\right)$, and $\arctan x =
\theta$ means $\tan\theta = x$ with $-\pi/2 < \theta < \pi/2$.
\end{definition}

Hence $\arctan(0) = 0$, $\arctan(1) = \pi/4$, and $\arctan(-1) = -\pi/4$.
None of these three intervals was chosen arbitrarily: each is a maximal
interval on which its function is monotonic and therefore one-to-one.
Restricting the domain further than necessary would discard values that
could otherwise be recovered, while enlarging it even slightly would
destroy invertibility by reintroducing repeated outputs.

\section{Composition Identities}

Inverse functions obey a pair of composition rules, and the two directions
of composition behave quite differently from one another.

\begin{theorem}
For every $x \in [-1,1]$, $\sin(\arcsin x) = x$.
\end{theorem}

\begin{proof}
By definition, $\arcsin x$ is the angle whose sine equals $x$; applying
sine to it simply recovers that value, $x$.
\end{proof}

The analogous statements $\cos(\arccos x) = x$ and $\tan(\arctan x) = x$
hold for the same reason. The reverse composition, however, is far more
delicate.

\begin{theorem}
$\arcsin(\sin\theta) = \theta$ holds only when $-\pi/2 \le \theta \le
\pi/2$.
\end{theorem}

\begin{proof}
By definition, $\arcsin$ must return a value lying in its restricted range
$\left[-\pi/2,\pi/2\right]$. If $\theta$ already lies in that range, the
returned value is $\theta$ itself. If $\theta$ lies outside that range,
$\arcsin(\sin\theta)$ instead returns whichever angle inside
$\left[-\pi/2,\pi/2\right]$ shares the same sine value as $\theta$, which is
a different number from $\theta$ whenever $\theta$ itself lies outside that
interval.
\end{proof}

This asymmetry between the two composition directions is one of the most
common sources of error in trigonometry, precisely because the forward
composition $\sin(\arcsin x) = x$ always holds while the reverse
composition does not, and it is worth working through carefully before
relying on either.

\begin{example}
Evaluate $\arcsin\left(\sin\dfrac{3\pi}{4}\right)$. First, $\sin(3\pi/4) =
\sqrt2/2$. The angle in $\left[-\pi/2,\pi/2\right]$ whose sine equals
$\sqrt2/2$ is $\pi/4$, so
\[
\arcsin\left(\sin\frac{3\pi}{4}\right) = \frac{\pi}{4},
\]
not $3\pi/4$: the inverse function is required to return a value inside its
restricted range, regardless of what angle originally produced the sine
value inside the parentheses.
\end{example}

\begin{algorithmproc}{How to Evaluate an Inverse Trigonometric Expression}
Identify which inverse function is involved and recall its restricted
output range. If the expression contains an inner trigonometric function
applied to an angle, compute that inner value first, then find the unique
angle in the restricted range that produces it. Before concluding, check
explicitly whether a composition identity is being applied, and never
assume $\arcsin(\sin\theta) = \theta$, or the analogous statement for cosine
or tangent, without first verifying that $\theta$ already lies in the
permitted range.
\end{algorithmproc}

\begin{practicenow}
Evaluate $\arccos\left(\cos\dfrac{5\pi}{4}\right)$.
\end{practicenow}

\section{Finding Angles from Ratios}

Inverse trigonometric functions allow an unknown angle to be solved for
directly rather than estimated. If $\tan\theta = 7/4$, applying $\arctan$
to both sides gives $\theta = \arctan(7/4)$, which a calculator evaluates
to approximately $1.052$ radians, or $60.3^\circ$. This formalizes, with a
precise definition behind it, the process first used informally in
Chapter~3 to recover an unknown angle from a known ratio of sides.

\section{Right-Triangle Interpretation}

An inverse trigonometric expression can also be read geometrically. Suppose
$\theta = \arcsin(3/5)$, so that $\sin\theta = 3/5$. Interpreting this
ratio as $\text{opposite}/\text{hypotenuse} = 3/5$ suggests a right
triangle with sides $3$, $4$, $5$, from which the remaining ratios follow
at once: $\cos\theta = 4/5$ and $\tan\theta = 3/4$. This construction turns
a single inverse trigonometric expression into a complete right triangle,
from which every other trigonometric quantity associated with $\theta$ can
be read off directly, without ever computing $\theta$ itself numerically.

\section{Applications}

Inverse trigonometric functions arise whenever an angle must be recovered
from geometric information already in hand: determining a ship's bearing
from coordinate data, computing an angle of elevation, analyzing forces in
engineering, recovering a direction from vector components, and solving
problems in navigation and surveying all rely on them. Together with the
trigonometric functions themselves, they complete a cycle running from
angle to ratio and back to angle again, closing the loop that Chapter~3
opened.

\section{Looking Ahead}

Inverse trigonometric functions solve equations of the simple forms
$\sin\theta = a$, $\cos\theta = a$, and $\tan\theta = a$. Many equations
that arise in practice are considerably more complicated, involving
products, powers, multiple angles, or several trigonometric functions at
once, and such equations cannot be solved systematically until a richer
collection of identities is available. The next several chapters develop
the geometry of rotation and the algebra of trigonometric identities that
this richer toolkit requires, after which Chapter~11 returns to the problem
of solving general trigonometric equations with the full apparatus in
place.

\section{Exercises}
\begin{exercise}
Evaluate $\arcsin\left(\dfrac{\sqrt3}{2}\right)$.
\end{exercise}

\begin{exercise}
Evaluate $\arccos\left(-\dfrac12\right)$.
\end{exercise}

\begin{exercise}
Evaluate $\arctan(1)$.
\end{exercise}

\begin{exercise}
Evaluate $\arcsin\left(\sin\dfrac{5\pi}{6}\right)$.
\end{exercise}

\begin{exercise}
Evaluate $\arccos\left(\cos\dfrac{7\pi}{6}\right)$.
\end{exercise}

\begin{exercise}
Find $\theta = \arctan\left(\dfrac53\right)$ to the nearest tenth of a
degree.
\end{exercise}

\begin{exercise}
If $\theta = \arcsin\left(\dfrac{8}{17}\right)$, find $\cos\theta$ and
$\tan\theta$.
\end{exercise}

\begin{exercise}
Explain why $\sin(\arcsin x) = x$ holds for every $x \in [-1,1]$, but
$\arcsin(\sin\theta) = \theta$ does not hold for every $\theta$.
\end{exercise}

\begin{exercise}
Construct a right triangle corresponding to $\theta = \arccos\left(\dfrac{12}{13}\right)$,
and use it to determine $\sin\theta$ and $\tan\theta$.
\end{exercise}

\begin{exercise}
A surveyor measures the ratio $\text{height}/\text{horizontal distance} =
0.65$. Find the angle of elevation.
\end{exercise}
